\subsection{The class-three lifting image of locally matching class-two sections}

\begin{statement}{A.201}{Proved}
{Every locally point-or-cusp-matching class-two section lifts continuously to class three}
Let \(U=\mathbf P^1_\Q\setminus\{0,1,\infty\}\),
\(\Pi=\pi_1(U_{\overline\Q})\), and use the closed lower
central series \(\Gamma_i\Pi\).  Put
\[
 E_2=E_U/\Gamma_3\Pi,\qquad E_3=E_U/\Gamma_4\Pi.
\tag{A.201.1}
\]
Let \(\mathscr S^{\rm loc}_2\) be the geometric conjugacy
classes of continuous sections of \(E_2\to G_\Q\) whose
restriction at each finite prime is the image, up to geometric
conjugacy, of an actual full point or cusp section of \(U\).
Then the actual projection of section sets satisfies
\[
 \mathscr S^{\rm loc}_2\ \subseteq\
 \operatorname{im}\bigl(\operatorname{Sec}(E_3/G_\Q)
                  \longrightarrow\operatorname{Sec}(E_2/G_\Q)\bigr).
\tag{A.201.2}
\]
For every representative in this domain, both continuous
weight-three obstruction classes of A.43 are zero, including
their coefficient-power derived-limit parts.  There is no
finite-support restriction on its Kummer valuations.
The conclusion gives a global class-three lift of the specified
class-two section.  The requirement that a lift also match the
chosen local point or cusp sections is not part of (A.201.2).
No local section-classification theorem is used: the domain
explicitly retains the actual local comparison maps.
\end{statement}

\begin{proof}
\emph{The original quotient and a fixed continuous lift of sets.}
Fix the rational tangent \(\overrightarrow{01}\), and the
meridians \(x,y\) of A.38 and A.43.  For a coefficient prime
\(\ell\), write \(N_{i,\ell}\) for the geometric pro-\(\ell\)
class-\(i\) quotient, and \(E_{i,\ell}\) for its actual
arithmetic extension.  The group maps are the characteristic
pro-\(\ell\) and lower-central quotient maps of \(E_U\).
With \(C=[x,y]\), \(D=[C,x]\), and \(F=[C,y]\), A.43
gives the exact sequence of continuous arithmetic groups
\[
 1\longrightarrow Z_\ell
 \longrightarrow E_{3,\ell}\xrightarrow{p}E_{2,\ell}
 \longrightarrow1,\qquad
 Z_\ell=\Z_\ell(3)D\oplus\Z_\ell(3)F.
\tag{A.201.3}
\]
Here the symbol \(Z_\ell\) in this proof denotes this
rank-two kernel.  Its geometric conjugation action is trivial;
the arithmetic action is multiplication by \(\chi_\ell^3\).
The literal Hall coordinates supply the continuous map of sets
\[
 L:N_{2,\ell}\longrightarrow N_{3,\ell},\qquad
        (a,b,h)\longmapsto y^bx^aC^h,
                  \qquad pL=\operatorname{id}_{N_{2,\ell}}.
\tag{A.201.4}
\]
All powers are continuous procyclic powers.  The map of sets
on arithmetic extensions is \((n,g)\mapsto(L(n),g)\), using
the same tangential splitting on both sides.

Choose an actual section \(s\) representing a class in
\(\mathscr S^{\rm loc}_2\), and project it to \(E_{2,\ell}\).
Its geometric cocycle is \(n(g)\in N_{2,\ell}\).  The
central defect of this fixed lift is the continuous cocycle
\[
 R(g,h)=L(n(g))\,g(L(n(h)))\,L(n(gh))^{-1}\in Z_\ell.
\tag{A.201.5}
\]
It is the raw defect in A.43.10; its two cohomology classes
are exactly \(o_D(s),o_F(s)\).  A continuous central cochain
\(z\) with \(\mathrm dz=-R\) supplies an actual section
\(g\mapsto(z(g)L(n(g)),g)\) of \(E_{3,\ell}\).

\emph{The exact inertia images of the local sections.}
Fix \(q\ne\ell\).  The pointed pro-\(\ell\) arithmetic
action on the tripod is unramified at \(q\).  This is the
prime-to-residue-characteristic specialization of the smooth
pair \((\mathbf P^1_{\Z_q},\{0,1,\infty\})\), with its
unit tangent.  More explicitly, the comparison of
\cite[Expos\'e XIII, Corollary 2.9 and \S2.10,
pp.288--290]{AgentSGA1} applies to finite \(\ell\)-groups,
since \(q\ne\ell\).  The tangent fibre functor is defined
using the compatible roots of the parameter and leading
coefficient one.  These roots and all prime-to-\(q\) roots
of unity are unchanged by the constant inertia group.
The comparison therefore makes its action trivial on each
pointed finite geometric quotient, and hence on
\(N_{3,\ell}\).  An integral ordinary point reducing away
from the boundary extends to this smooth model; its pulled-back
finite covers are unramified.  Its geometric inertia image
is consequently the identity.

For an ordinary point reducing to a boundary point, the maps
\(\operatorname{Spec}\Z_q[[T]][1/(qT)]\to U_{\Q_q}\)
and \(\operatorname{Spec}\Q_q((T))\to
\operatorname{Spec}\Z_q[[T]][1/(qT)]\) of A.49 identify
the full geometrically pro-\(\ell\) section with the
corresponding cusp section.  The literal parameters are
\(T=t,1-t,1/t\), at \(0,1,\infty\), respectively.
The tame normalization and constant-extension maps in
A.49.3--6 prove that the geometric coordinate is the
Kummer character of this same parameter.  Equivalently,
the normal-crossings Kummer cover in
\cite[Expos\'e XIII, Propositions 5.1--5.2,
pp.314--315]{AgentSGA1} kills its prime-to-\(q\) ramification;
evaluation at the point sends \(T^{1/\ell^a}\) to the
corresponding root of its value.  This specifies both the
arithmetic map and its geometric conjugation ambiguity.
For an actual cusp section, its completed tangent Kummer
parameter gives the same description directly.

Let \(T_q\) be a tame \(\ell\)-inertia generator normalized
by \(\kappa_\ell(q)(T_q)=1\).  Thus every local section just
described has, in \(N_{3,\ell}\), an inertia value of the form
\[
             u\,c^{m_q}u^{-1},\qquad
 c\in\{x,y,(xy)^{-1}\},\quad u\in N_{3,\ell},
                              \quad m_q\in\Z_\ell,
\tag{A.201.6}
\]
or the identity.  At ordinary annular points the exponent is
the ordinary valuation of the displayed parameter.  At a cusp
it is the valuation component of its completed tangent class.
The discarded wild and prime-to-\(\ell\) inertia acts trivially
and maps trivially to the pro-\(\ell\) geometric group, as
in A.46.  All calculations therefore take place on the same
tame quotient used by its strict residue map.

By the definition of \(\mathscr S^{\rm loc}_2\), a local
geometric conjugator aligns the actual full point or cusp
section with \(s\) in \(E_{2,\ell}\).  Lift that conjugator
through \(N_{3,\ell}\to N_{2,\ell}\), and conjugate the
full local section before taking its class-three quotient.
This supplies an actual local section \(\sigma_q\) satisfying
\(p\sigma_q=s|_{G_{\Q_q}}\) literally.  Its inertia still
has the form (A.201.6), because the background inertia action
is trivial.  The two global abelian inertia coordinates satisfy
\[
 (a(T_q),b(T_q))=(m_q,0),\quad(0,m_q),\quad
                     (-m_q,-m_q),\quad\hbox{or }(0,0).
\tag{A.201.7}
\]
In the last case choose the identity and \(m_q=0\).
Every finite reduction of the two global Kummer cocycles is
unramified outside finitely many primes.  Consequently the
exponents selected in (A.201.7) satisfy
\[
                         m_q\longrightarrow0\quad\hbox{in }\Z_\ell
\tag{A.201.8}
\]
as \(q\) leaves finite sets.  This applies equally to the
completed tangent exponents.

\emph{A uniform integral estimate in the actual Hall coordinates.}
Define the continuous central-coordinate map
\[
 \mathcal F_c(u,m)=u c^m u^{-1}
                  L\bigl(p(u c^m u^{-1})\bigr)^{-1}
                         \in Z_\ell.
\tag{A.201.9}
\]
It is zero at \(m=0\), uniformly on the compact group
\(N_{3,\ell}\).  The following explicit estimate retains
the small coefficient primes.  The central part of \(u\)
cancels from its conjugation, so write its remaining part as
\(y^sx^rC^h\).  With coordinates on \((D,F)\), the three
maps in (A.201.9) are
\[
\begin{aligned}
 \mathcal F_{x,D}&=m(2h-ms-s)/2,
 &\mathcal F_{x,F}&=ms(s-1)/2,\\
 \mathcal F_{y,D}&=-mr(r-1)/2,
 &\mathcal F_{y,F}&=m(2h+mr-2rs+r)/2.
\end{aligned}
\tag{A.201.10}
\]
For \(c=(xy)^{-1}\), they are
\[
\begin{aligned}
 6\mathcal F_{c,D}
    &=m(-6h-2m^2+3mr-3ms+3m+3r^2+3s-1),\\
 6\mathcal F_{c,F}
    &=m(-6h-m^2+3mr-3ms+3m+6rs-3r-3s^2+6s-2).
\end{aligned}
\tag{A.201.11}
\]
These are identities of integral group words.  They follow by
substituting \(x=1+X,y=1+Y\) in the noncommutative Magnus
ring truncated after degree three, using
\((1+X)^r=1+rX+\binom r2X^2+\binom r3X^3\), and the
same collecting convention as A.43.7--10.  The negative
\(X^2Y\) coefficient and the \(XY^2\) coefficient extract
the \(D,F\) coordinates after removing the Hall lift.
The exact certificate checks those identities over the rational
polynomial ring, including that every coefficient of
\(6\mathcal F_c/m\) is an integer.  Integral powers first
give the identities; continuous procyclic powers extend them
to \(r,s,h,m\in\Z_\ell\).  Their values are integral
because they are the central coordinates of (A.201.9), even
when the printed denominators are not units.  In particular
each coordinate obeys
\[
 v_\ell(\mathcal F_c(u,m))\ge
                  v_\ell(m)-v_\ell(6).
\tag{A.201.12}
\]
This proves the required convergence uniformly in every choice
of local conjugator, with an explicit loss of at most one
power at two or three.

\emph{The strict continuous residue, with its actual denominator.}
Write the aligned local section uniquely as
\(\sigma_q(g)=(z_q(g)L(n(g)),g)\), with a continuous
central cochain \(z_q:G_{\Q_q}\to Z_\ell\).
Its homomorphism identity gives
\(\mathrm dz_q=-R|_{G_{\Q_q}}\).  Its inertia value is
exactly
\[
                    z_q(T_q)=\mathcal F_c(u_q,m_q),
\tag{A.201.13}
\]
with zero in the integral ordinary case.  Equations
(A.201.8) and (A.201.12) imply that this family tends to zero.

For clarity, the strict residue map from A.27 and A.46, for
weight-three coefficients, sends a normalized two-cochain
\(V\) to
\[
 \psi_q^2(V)=q^{-1}\left(
 V(\Phi_q,T_q)+V(\Phi_qT_q,\Phi_q^{-1})
 -V(\Phi_q,\Phi_q^{-1})
 -\sum_{j=1}^{q-1}V(T_q^j,T_q)\right).
\tag{A.201.14}
\]
Here \(\Phi_qT_q\Phi_q^{-1}=T_q^q\), and the same
fixed choices are used at every coefficient level.  The exact
telescoping identity is
\(\psi_q^2(\mathrm dz)=(q^2-1)z(T_q)\).
Consequently the residue family of the original global defect
in (A.201.5) is literally
\[
          \psi_q^2(R)=-(q^2-1)z_q(T_q).
\tag{A.201.15}
\]
This identifies its primitive before quotienting any local or
global cohomology group.

Put \(R_\ell^0=\{(w_q)_{q\ne\ell}:w_q\in\Z_\ell,
w_q\to0\}\), with cofinite convergence.  A.46.2 gives the
actual continuous isomorphism
\[
 H^2_{\rm cts}(G_\Q,\Z_\ell(3))
 \simeq R_\ell^0/
       \{((q^2-1)w_q)_q:(w_q)_q\in R_\ell^0\}.
\tag{A.201.16}
\]
Its proof retains the termwise direct sums, the inverse-limit
cochain comparison, and the literal image in the denominator.
The unramified term is zero at odd \(\ell\) by
\cite[VI, Theorem 8.2, with \(n=4\)]{Weibel2013} and
\cite[Theorem 0.1]{BIR_Rognes2000}; at two A.46 proves its
vanishing by restriction to \(\Q(i)\) and reduction at the
real place for the sign module \(\Z_2(3)\).  Thus this
comparison includes every coefficient prime and the real
cohomological term.
By (A.201.13), the two-coordinate family \((-z_q(T_q))_q\)
belongs to \((R_\ell^0)^2\).  Equation (A.201.15) therefore
places the full global defect in the literal denominator of
(A.201.16).  Both continuous obstruction classes vanish.

Finally the geometric nilpotent class-three quotient is the
product of its pro-\(\ell\) quotients: each finite quotient
is nilpotent and is the direct product of its Sylow subgroups,
and the actual inverse-limit maps give this product decomposition.
Equivalently the five Hall coordinates of A.43 decompose
coefficient by coefficient.  Continuous cochains with values
in \(\widehat\Z(3)^2=\prod_\ell\Z_\ell(3)^2\) form the
product of their coefficient complexes.  Choose a continuous
primitive of \(-R\) in every factor.  Their product is
continuous and supplies the section of \(E_3\) in (A.201.2).
A change of representative of \(s\) is handled by lifting its
geometric conjugator through \(\Pi/\Gamma_4\Pi\to
\Pi/\Gamma_3\Pi\), so the claimed inclusion is indeed on
the stated geometric conjugacy classes.

The finite symbolic identities (A.201.10)--(A.201.12) are
verified by
\path{code/selmer_class_two_to_three_certificate.py}, with
retained output and exact data.  The continuous primitive is
proved by (A.201.13)--(A.201.16).  This result specifies a
lifting image from the locally matching class-two domain;
the comparison of arbitrary ordinary Kummer pairs with that
domain, and further locally matching lifts, remain separate
unresolved steps toward A1 and A5.
\end{proof}
